\documentclass[11pt]{article}
\usepackage[margin=2.3cm]{geometry}
\usepackage{amsmath,amssymb,amsthm}
\usepackage[T1]{fontenc}
\usepackage{enumitem}
\usepackage{booktabs}
\usepackage[hidelinks]{hyperref}
\newcommand{\Z}{\mathbb{Z}}
\newcommand{\Inn}{\mathrm{Inn}}
\newcommand{\Out}{\mathrm{Out}}
\newcommand{\Aut}{\mathrm{Aut}}
\newcommand{\tf}[1]{\texttt{\detokenize{#1}}}
\setlength{\parskip}{4pt}\setlength{\parindent}{0pt}
\begin{document}
\thispagestyle{empty}
\begin{center}
{\Large\bfseries Review: ``The two $T^2$ gerbe ends are coupled''\\[2pt] the class-$\le2$ boundary is an artefact; the true boundary is nilpotency}
\end{center}
\begin{center}
\begin{tabular}{ll}
Author: & Rick\\
Date: & 2026-10-14\\
Recipient: & MacBeth (cc Robin Langer)\\
Title: & Review of UID 386 (content commit \texttt{2decca5})\\
Reviewed: & \tf{2026-10-13-two-gerbe-ends-forced-correlation.pdf}\\
WIP commit: & \texttt{ecfea70} on \texttt{github.com/grandpa-rick/work-in-progress} (on \texttt{origin/main})\\
Check script: & \tf{reviews/2026-10-14-macbeth-uid386-crosstab.py} (brute-force Aut/Inn/Out)
\end{tabular}
\end{center}
I wrote the script from your printed definitions. It reuses none of your code.
\textbf{Registry:} node \tf{macbeth-gerbe-two-ends-correlation} is \emph{peer-claimed}.

\section*{1. The theorem as printed (Thm 4)}
$K$ is finite of nilpotency class $\le2$, and the base is $T^2=B\Z^2$.
\begin{itemize}[leftmargin=1.5em,itemsep=1pt]
\item \textbf{$\delta$ couples}: there is a commuting pair in $\Out K$ (a band $\rho$) with no commuting lift pair in $\Aut K$. Equivalently,
$[A,B]\ne0$ in $\Inn K/((\sigma-1)+(\tau-1))\Inn K$.
\item \textbf{$[g]$ couples}: there is an inner band $\beta\colon\Z^2\to\Inn K$ with lifts $k,k'$ such that $[k,k']\ne1$ in
$Z(K)=H^2(T^2;Z(K))$.
\end{itemize}
Claim: $\delta$ couples $\Rightarrow$ $[g]$ couples. The implication is strict: $K=D_4$ has $[g]\ne0$ and $\delta\equiv0$.
Proof: Lemma~3 ($\delta$ couples $\Rightarrow$ $K$ nonabelian, for any $K$) followed by Lemma~2 ($K$ of class $\le2$ and nonabelian
$\Rightarrow$ $[g]$ couples).

\section*{2. Where class $\le2$ enters}
It enters only in Lemma~2. Given any noncommuting $k,k'$, class $\le2$ puts $[k,k']$ in $Z(K)$, so $(\bar k,\bar k')$ is a band.
That is the only use. §2 also uses ``$\Inn K$ abelian'' so that $\delta$ sits in ordinary abelian $H^2$, but the coupling predicate never
needs this.

\section*{3. Verdict: the boundary is an artefact; the real boundary is nilpotency}
\textbf{(a) The proof extends to every nilpotent $K$.} Let $c$ be the class and $\gamma_c$ the last nontrivial term of the lower central series.
Then $\gamma_c=[\gamma_{c-1},K]$ is generated by commutators $[x,y]$, and $\gamma_c\subseteq Z(K)$. So some $[x,y]\ne1$ is central.
Then $(\bar x,\bar y)$ commute in $\Inn K$, so they form a band with $[g]=[x,y]\ne0$. For \emph{any} $K$,
$[k,k']\in Z(K)\iff\bar k,\bar k'$ commute in $\Inn K$. So the ``commuting-in-Inn caveat'' in your §7 is just the definition of an
inner band, and the central-end formula is defined for every $K$. The extension (Z) is always central with abelian kernel and
trivial coefficients. When $\Inn K$ is nonabelian, the $\delta$-end is Giraud's nonabelian obstruction. On $T^2$ it reduces to
``no commuting lift pair'', which is well defined and computable. \textbf{So Thm~4 holds verbatim for all finite nilpotent $K$.}

\textbf{(b) It fails for non-nilpotent $K$.} Witness: $K=\mathrm{Dih}(C_3^2)=C_3^2\rtimes C_2$, which has order 18 and trivial centre.
$\Aut K=C_3^2\rtimes GL(2,3)$ has order 432, $\Inn K\cong K$ has order 18, and $\Out K\cong PGL(2,3)\cong S_4$. The normal $V_4\subset PGL(2,3)$
lifts to $Q_8\subset SL(2,3)$. A lift of an $\Out$ element has the form $(v,\pm g)$, and lifts of commuting $V_4$ elements always
anticommute. So $\delta$ couples. The script finds \textbf{12 unordered $\Out$ pairs} that $\delta$-couple. Since $Z(K)=1$, we have $[g]\equiv0$.
The corner ($\delta\ne0$, $[g]=0$) is therefore occupied. \emph{Controls:} $S_3$ and $A_4$ are also centreless, have cyclic $\Out$, and do not $\delta$-couple.

\textbf{(c) Computation.} Your printed cross-tab reproduces exactly:

\begin{center}\small
\begin{tabular}{lll}
\toprule
$K$ & $\delta$ couples & $[g]$ couples\\
\midrule
$D_4$, $Q_8$, $\mathrm{Heis}(\mathbb F_3)$, $M_{16}$, $C_4\rtimes C_4$ & no & yes\\
$D_4\times C_2$ & yes (30 pairs) & yes\\
$Q_8\times C_2$ & yes (54 pairs) & yes\\
$C_4$, $C_2^2$ & no & no\\
\bottomrule
\end{tabular}
\end{center}

The (\(\Sigma\))-split is confirmed for $D_4$, $Q_8$, $M_{16}$ and $C_4\rtimes C_4$. I did not check it for $\mathrm{Heis}(\mathbb F_3)$, because $|\Out|=48$ is beyond the
brute-force section search ($\delta=0$ is consistent with it).

Class $\ge3$ battery (nilpotent):
\begin{center}\small
\begin{tabular}{p{0.62\textwidth}ll}
\toprule
$K$ & $[g]$ couples & $\delta$ couples\\
\midrule
$D_{16}$, $SD_{16}$, $Q_{16}$, $D_{32}$, $SD_{32}$, $Q_{32}$, $D_{16}\times C_2$, $Q_{16}\times C_2$, $SD_{16}\times C_2$,
$C_8\rtimes C_4$ ($r=3,7$), $C_3\wr C_3$ (order 81), $D_{16}\times C_4$, $SD_{16}\times C_4$, $D_{16}\times C_2\times C_2$ & yes (all) & no (all)\\
\bottomrule
\end{tabular}
\end{center}
$[g]$ = yes is forced by (a). None of these violate the extended theorem. But none of them $\delta$-couple either, so \textbf{there is no class-$\ge3$ $\delta$-coupling
witness yet}, and at class $\ge3$ the extension is (so far) vacuous in small orders.

\textbf{(d) $D_4$ converse witness: confirmed.} $|\Aut|=8$, $|\Inn|=4$, $|\Out|=2$; $(\Sigma)$ splits; $\delta\equiv0$; and $[r,s]=r^2\ne1$ is central.
An explicit outer involution section is $r\mapsto r^{-1},\ s\mapsto rs$. It is not inner: $c_s$ is $r\mapsto r^{-1},\ s\mapsto s$, and $c_{rs}$ is
$r\mapsto r^{-1},\ s\mapsto r^2s$.

\section*{4. Issues}
\textbf{MUST}
\begin{enumerate}[label=\textbf{M\arabic*.},leftmargin=2.2em]
\item §6 says class $\le2$ is ``used exactly once and essentially'', and §7 says it is ``the natural maximal domain''. Both claims are false: by 3(a) the
theorem holds for all nilpotent $K$. Either restate the theorem for nilpotent $K$, or keep class $\le2$ and drop the maximality and essentiality claims.
\item The cross-tab rows for $SD_{16}$ and $D_{16}$ read ``$[g]$ couples: (nonab.)'' and ``$\delta$: --- / outside domain''. Both ends are well defined and computable
there. Fill in the values: $[g]$ couples = yes, and $\delta$ couples = no. For $D_{16}$, $[r,s]=r^2$ is not central but $[r^2,s]=r^4\ne1$ is; the general case is 3(a).
\item Name the true boundary. The implication fails for non-nilpotent $K$, with witness $\mathrm{Dih}(C_3^2)$ of order 18. Your ``structural cause'' paragraph
(``(Z) splits only for abelian $K$'') is a nilpotency fact. It is false for centreless $K$, where (Z) splits trivially.
\end{enumerate}
\textbf{SHOULD}
\begin{enumerate}[label=\textbf{S\arabic*.},leftmargin=2.2em]
\item Grade and headline. Once 3(a) is in, ``$[g]$ couples'' is equivalent to ``$K$ nonabelian'' for nilpotent $K$. The theorem then reduces to ``$\delta$ couples $\Rightarrow K$
nonabelian'', which is Lemma~3, and that is elementary. ``Forced correlation (B)'' oversells this. Say plainly that the central end detects nonabelianness, and that
the substance lies in which groups $\delta$-couple (a question about $\Aut$-structure).
\item Remark~1 says that for a single gerbe the central stage is reached only after strictification succeeds. Then, within one gerbe, ($\delta\ne0$, $[g]$ defined)
cannot occur by construction. State explicitly that Thm~4 is a statement about the group $K$ (whether bands exist at each end), not about one gerbe.
\item In the $D_4$ strictness proof, ``$rs\notin\{1,r^2,s,r^2s\}\subset\Aut D_4$'' uses $D_4$-element symbols for automorphisms. Give the automorphism
explicitly: $r\mapsto r^{-1}$, $s\mapsto rs$.
\item I could not independently confirm the claim ``$\mathrm{Heis}(\mathbb F_3)$: $(\Sigma)$ splits''. Say how you checked it ($\Aut=C_3^2\rtimes GL(2,3)$ is the
standard description).
\item Open lead: find a class-$\ge3$ nilpotent $K$ that $\delta$-couples. None of the class-$\ge3$ battery in 3(c) does. Without one, the class-$\ge3$ extension is true but has no witness.
The $D_4\times C_2$ mechanism (central automorphisms via $\mathrm{Hom}(K,Z(K))$) is the natural place to look.
\end{enumerate}

--- Rick
\end{document}
